\section{Results}

\subsubsection{Assurance Coverage}

Of the 50 target functions, 5 (10\%) are \emph{Verified}---Nagini formally discharged the equivalence proof between early and final assertions; 30 (60\%) are \emph{Validated}---they passed both CrossHair symbolic execution and Hypothesis fuzz testing but were unsupported or inconclusive under Nagini; and 11 (22\%) are \emph{Rejected}---at least one stage produced a counterexample (Detailed results in Table \ref{tab:program-results}). Notably, the Verified group is dominated by pure-arithmetic functions with minimal control flow, whereas functions that rely on randomness, floating-point operations, or complex stdlib calls appear exclusively in the Rejected category.


\subsubsection{Fail-Fast Security Impact}

Moving assertions earlier in the program yields tangible security and efficiency benefits. Across our benchmark, early assertion placement saved an average of 28.6 Python bytecode instructions per function (median 20), with over 80\% of failed inputs halted within the first 10 instructions (Figure \ref{fig:bytecode-load}). This reduction in executed instructions before error detection directly shrinks the window for side-effects, resource exhaustion, or side-channel leakage. The greatest savings were observed in pure arithmetic and simple loop functions, which typically halt in under 15 instructions, while string manipulation and floating-point operations required more computation before reaching the assertion. These results highlight the practical value of fail-fast enforcement, especially in security-sensitive or resource-constrained contexts.

\subsubsection{Complexity vs. Verifiability}

Verification success rates are strongly tied to program complexity. As shown in Figure \ref{fig:diff-scores-correlation}, functions with overall difficulty scores below 2.5 passed CrossHair over 80\% of the time and Nagini over 40\%, while those with scores above 4.0 rarely succeeded under either tool. The average difficulty of Verified functions was 2.0, compared to 2.7 for Rejected ones (Table \ref{tab:difficulty-sub-scores}).

A closer look at the rejected functions (Table \ref{tab:program-results}) reveals several recurring patterns. Nearly all rejected cases feature at least one of: deep or nested control flow (e.g., multiple loops or branches), floating-point arithmetic, or use of randomness and non-deterministic library calls. For example, functions like \texttt{binary\_search\_iterations}, \texttt{matrix\_determinant}, and \texttt{polygon\_area\_calculator} combine high control-flow depth with complex data manipulations, while \texttt{convert\_temperature} and \texttt{calculate\_discount} rely on floating-point comparisons that are notoriously difficult for SMT-based tools like CrossHair and Nagini. Randomness-heavy functions (e.g., \texttt{random\_mod\_calculator}) are systematically rejected, as symbolic and static verifiers cannot reason about non-deterministic outcomes. 

Additionally, list comprehensions, dynamic data structures, and external library calls (such as math or datetime) are overrepresented among failures. These features increase the search space and introduce behaviors that are hard to model or prove correct in a general way. This analysis underscores that, while our pipeline is robust for simple, linear code, it currently struggles with the kinds of dynamic, data-rich, or numerically sensitive logic that typifies real-world Python programs. Addressing these limitations is a key direction for future work in automated formal verification for security.

\subsubsection{Diagnostic Richness of Failures}

We were able to analyze 18 separate failure cases across fuzzing, CrossHair, and Nagini, where we investigated syntax errors, conditional failures, counterexamples, and proof errors. Across all failed cases, the mean Failure Explanation Quality (FEQ) score is 0.53 (Table \ref{tab:feq-scores}), indicating that most counterexamples and error messages are specific, actionable, and reasonably context-rich. Failures involving floating-point operations or complex data manipulations tend to yield the most detailed explanations (e.g., \texttt{convert\_temperature}, \texttt{digit\_sum\_processor}), with FEQ scores as high as 0.74. In contrast, failures in randomness-heavy or highly dynamic code are less informative, as the tools struggle to generate concrete counterexamples or detailed failure points in program execution.

Not all failures, however, reflect true logical inequivalence of early and final assertions. A subset of five cases---marked as 'inconclusive (false positive)' in Table \ref{tab:program-results} and detailed in Table \ref{tab:feq-scores}---arise from tool limitations, unsupported language features, or type errors rather than genuine assertion mismatches. For example, \texttt{isbn\_validator} and \texttt{matrix\_determinant} both failed due to dynamic data manipulations and unsupported operations, while \texttt{date\_difference\_calculator} and \texttt{day\_of\_week\_calculator} triggered false positives because of reliance on Python's datetime library, which is not modeled by the verification tools. \texttt{convert\_temperature} produced a spurious failure due to floating-point precision issues. In our analysis, we distinguish these false positives from genuine counterexamples, ensuring that only true logical failures are counted as rejections. This distinction is critical for accurately assessing both the strengths and the current boundaries of automated formal verification in practice.

Notably, functions that fail both CrossHair and Nagini typically receive the highest FEQ scores, benefiting from the complementary strengths of symbolic and static analysis. This multi-tool approach not only increases the likelihood of catching subtle errors but also improves the quality of diagnostic feedback, helping developers quickly identify and address the root causes of assertion inequivalence.

\subsubsection{Residual Attack Surface and Tool Support Gaps} 

Despite the strengths of our pipeline, a significant fraction of real-world Python code remains outside the reach of current formal and symbolic verification tools. As detailed in Table \ref{tab:nagini-unsupported-programs}, the most common unsupported features are floating-point arithmetic (80\% of unsupported cases), list comprehensions (65\%), randomness (40\%), and datetime or external library calls (30\%). These features either introduce non-determinism, require complex heap modeling, or depend on external semantics that are not captured by the verification engines.

Representative examples include \texttt{date\_difference\_calculator} and \texttt{day\_of\_week\_calculator} (unsupported due to datetime operations), \texttt{isbn\_validator} and \texttt{loop\_string\_hash} (dynamic comprehensions and string/list indexing), \texttt{matrix\_determinant} (modular arithmetic and matrix ops), and \texttt{random\_mod\_calculator} (randomness). Many mathematical and data-rich programs, such as \texttt{factorial\_root\_calculator}, \texttt{polygon\_area\_calculator}, and \texttt{mean\_absolute\_deviation}, are also unsupported due to their use of advanced math functions, rounding, or complex control flow.

As a result, these features define the current "blind spots" in automated assertion verification, and represent key targets for future tool development. Addressing these gaps will be essential for scaling formal security guarantees to the full diversity of Python programs encountered in practice.


